\chapter{Instruction Summary}
\label{ch:instruction-summary}

\section{Overview}

The SIRCIS instruction set consists of 64 possible opcodes (6-bit opcode field), of which 50 are documented and
assigned; the remaining 14 are undocumented (see Appendix~\ref{appendix:undocumented}). The instructions are organized
into logical groups:

\begin{itemize}
	\item \textbf{ALU Instructions (\opcode{00}--\opcode{0E}, \opcode{20}--\opcode{2E}, \opcode{30}--\opcode{3E})}: Arithmetic, logical, and comparison operations
	\item \textbf{Memory Instructions (\opcode{10}--\opcode{17})}: Load and store operations
	\item \textbf{Address and Control Flow (\opcode{18}--\opcode{1F})}: Effective-address calculation, jumps, and calls
	\item \textbf{Coprocessor (\opcode{0F}, \opcode{3F})}: Coprocessor interface
\end{itemize}

\opcode{2F} is undocumented (see Appendix~\ref{appendix:undocumented}).

\section{Complete Instruction List}

\begin{table}[H]
	\centering
	\tiny
	\begin{tabular}{clllp{4cm}}
		\toprule
		\textbf{Opcode} & \textbf{Mnemonic} & \textbf{Format} & \textbf{Flags} & \textbf{Operation}                              \\
		\midrule
		\opcode{00}     & ADDI              & Immediate       & NZCV           & Add immediate                                   \\
		\opcode{01}     & ADCI              & Immediate       & NZCV           & Add immediate with carry                        \\
		\opcode{02}     & SUBI              & Immediate       & NZCV           & Subtract immediate                              \\
		\opcode{03}     & SBCI              & Immediate       & NZCV           & Subtract immediate with carry                   \\
		\opcode{04}     & ANDI              & Immediate       & NZ (C,V=0)     & AND immediate                                   \\
		\opcode{05}     & ORRI              & Immediate       & NZ (C,V=0)     & OR immediate                                    \\
		\opcode{06}     & XORI              & Immediate       & NZ (C,V=0)     & XOR immediate                                   \\
		\opcode{07}     & LOAD              & Immediate       & --             & Load immediate (move)                           \\
		\opcode{08}     & --                & Immediate       & --             & Undocumented                                    \\
		\opcode{09}     & --                & Immediate       & --             & Undocumented                                    \\
		\opcode{0A}     & CMPI              & Immediate       & NZCV           & Compare immediate                               \\
		\opcode{0B}     & --                & Immediate       & --             & Undocumented                                    \\
		\opcode{0C}     & TSAI              & Immediate       & NZ (C,V=0)     & Test AND immediate                              \\
		\opcode{0D}     & --                & Immediate       & --             & Undocumented                                    \\
		\opcode{0E}     & TSXI              & Immediate       & NZ (C,V=0)     & Test XOR immediate                              \\
		\opcode{0F}     & COPI              & Immediate       & --             & Coprocessor call immediate                      \\
		\midrule
		\opcode{10}     & STOR              & Indirect Imm    & --             & Store to (addr + imm)                           \\
		\opcode{11}     & STOR              & Indirect Reg    & --             & Store to (addr + reg)                           \\
		\opcode{12}     & STOR              & Pre-Dec         & --             & Store with pre-decrement                        \\
		\opcode{13}     & STOR              & Pre-Dec Reg     & --             & Store with pre-dec (reg)                        \\
		\opcode{14}     & LOAD              & Indirect Imm    & --             & Load from (addr + imm)                          \\
		\opcode{15}     & LOAD              & Indirect Reg    & --             & Load from (addr + reg)                          \\
		\opcode{16}     & LOAD              & Post-Inc        & --             & Load with post-increment                        \\
		\opcode{17}     & LOAD              & Post-Inc Reg    & --             & Load with post-inc (reg)                        \\
		\opcode{18}     & LDEA              & Indirect Imm    & --             & Load effective address                          \\
		\opcode{19}     & LDEA              & Indirect Reg    & --             & Load effective address (reg)                    \\
		\opcode{1A}     & LDEA              & Pre-Dec Imm     & --             & Load effective address, pre-dec                 \\
		\opcode{1B}     & LDEA              & Pre-Dec Reg     & --             & Load effective address, pre-dec (reg)           \\
		\opcode{1C}     & LDEL              & Indirect Imm    & --             & Load effective address and link                 \\
		\opcode{1D}     & LDEL              & Indirect Reg    & --             & Load effective address and link (reg)           \\
		\opcode{1E}     & LDEL              & Post-Inc Imm    & --             & Load effective address and link, post-inc       \\
		\opcode{1F}     & LDEL              & Post-Inc Reg    & --             & Load effective address and link, post-inc (reg) \\
		\midrule
		\opcode{20}     & ADDI              & Short Imm+Shift & NZCV           & Add short immediate                             \\
		\opcode{21}     & ADCI              & Short Imm+Shift & NZCV           & Add short imm with carry                        \\
		\opcode{22}     & SUBI              & Short Imm+Shift & NZCV           & Subtract short immediate                        \\
		\opcode{23}     & SBCI              & Short Imm+Shift & NZCV           & Subtract short imm with carry                   \\
		\opcode{24}     & ANDI              & Short Imm+Shift & NZ (C,V=0)     & AND short immediate                             \\
		\opcode{25}     & ORRI              & Short Imm+Shift & NZ (C,V=0)     & OR short immediate                              \\
		\opcode{26}     & XORI              & Short Imm+Shift & NZ (C,V=0)     & XOR short immediate                             \\
		\opcode{27}     & --                & Short Imm+Shift & --             & Undocumented                                    \\
		\opcode{28}     & --                & Short Imm+Shift & --             & Undocumented                                    \\
		\opcode{29}     & --                & Short Imm+Shift & --             & Undocumented                                    \\
		\opcode{2A}     & CMPI              & Short Imm+Shift & NZCV           & Compare short immediate                         \\
		\opcode{2B}     & --                & Short Imm+Shift & --             & Undocumented                                    \\
		\opcode{2C}     & TSAI              & Short Imm+Shift & NZ (C,V=0)     & Test AND short immediate                        \\
		\opcode{2D}     & --                & Short Imm+Shift & --             & Undocumented                                    \\
		\opcode{2E}     & TSXI              & Short Imm+Shift & NZ (C,V=0)     & Test XOR short immediate                        \\
		\opcode{2F}     & --                & Short Imm+Shift & --             & Undocumented                                    \\
		\midrule
		\opcode{30}     & ADDR              & Register        & NZCV           & Add register                                    \\
		\opcode{31}     & ADCR              & Register        & NZCV           & Add register with carry                         \\
		\opcode{32}     & SUBR              & Register        & NZCV           & Subtract register                               \\
		\opcode{33}     & SBCR              & Register        & NZCV           & Subtract register with carry                    \\
		\opcode{34}     & ANDR              & Register        & NZ (C,V=0)     & AND register                                    \\
		\opcode{35}     & ORRR              & Register        & NZ (C,V=0)     & OR register                                     \\
		\opcode{36}     & XORR              & Register        & NZ (C,V=0)     & XOR register                                    \\
		\opcode{37}     & LOAD              & Register        & --             & Load from register (move)                       \\
		\opcode{38}     & --                & Register        & --             & Undocumented                                    \\
		\opcode{39}     & --                & Register        & --             & Undocumented                                    \\
		\opcode{3A}     & CMPR              & Register        & NZCV           & Compare register                                \\
		\opcode{3B}     & --                & Register        & --             & Undocumented                                    \\
		\opcode{3C}     & TSAR              & Register        & NZ (C,V=0)     & Test AND register                               \\
		\opcode{3D}     & --                & Register        & --             & Undocumented                                    \\
		\opcode{3E}     & TSXR              & Register        & NZ (C,V=0)     & Test XOR register                               \\
		\opcode{3F}     & COPR              & Register        & --             & Coprocessor call register                       \\
		\bottomrule
	\end{tabular}
	\caption{Complete SIRCIS Instruction Set}
	\label{tab:complete-instruction-set}
\end{table}

See Table~\ref{tab:flag-effect-symbols} for the flag-effect symbol definitions. The Flags column shows \texttt{--} for
\mnemonic{LOAD} (\opcode{07}, \opcode{37}) only when the default \texttt{[N]} status update source is used; an
explicit status update override on \mnemonic{LOAD} leaves all four flags architecturally undefined. See
Chapter~\ref{ch:alu-instructions} for
details.

\section{Instruction Grouping Patterns}

\subsection{Opcode Organization}

The opcode space is systematically organized:

\begin{itemize}
	\item \textbf{\texttt{0x0\_}}: Immediate operand format
	\item \textbf{\texttt{0x1\_}}: Memory and control flow operations
	\item \textbf{\texttt{0x2\_}}: Short immediate with shift format
	\item \textbf{\texttt{0x3\_}}: Register operand format
\end{itemize}

\subsection{Save vs. Test Variants}

Many ALU instructions have two variants:

\begin{itemize}
	\item \textbf{+\opcode{00} to +\opcode{07}}: Save result to destination register
	\item \textbf{+\opcode{08} to +\opcode{0E}}: Test only (update flags, discard result)
	\item \textbf{+\opcode{0F}}: Coprocessor call (neither result nor flags are written)
\end{itemize}

Examples:
\begin{itemize}
	\item \opcode{04} = \mnemonic{ANDI} (save result), \opcode{0C} = \mnemonic{TSAI} (test only)
	\item \opcode{06} = \mnemonic{XORI} (save result), \opcode{0E} = \mnemonic{TSXI} (test only)
\end{itemize}

\section{Meta-Instructions}

Several commonly-used operations are implemented as assembler meta-instructions that assemble to specific opcodes:

\begin{table}[H]
	\centering
	\small
	\begin{tabular}{llp{7cm}}
		\toprule
		\textbf{Meta-Instruction} & \textbf{Assembles To}           & \textbf{Purpose}         \\
		\midrule
		\mnemonic{NOOP}           & \texttt{ADDI[N] r1, \#0}        & No operation             \\
		\mnemonic{RETS}           & \texttt{LDEA p, (\#0, l)}       & Return from subroutine   \\
		\mnemonic{SHFT}           & \texttt{ORRI[S] rD, \#0, shift} & Shift with status update \\
		\mnemonic{BRAN}           & \texttt{LDEA p, ...}            & PC-relative branch       \\
		\mnemonic{BRSR}           & \texttt{LDEL p, ...}            & PC-relative call         \\
		\mnemonic{LJMP}           & \texttt{LDEA p, ...}            & Long jump                \\
		\mnemonic{LJSR}           & \texttt{LDEL p, ...}            & Long call                \\
		\mnemonic{EXCP}           & \texttt{COPI \#0x11vv}          & User exception           \\
		\mnemonic{WAIT}           & \texttt{COPI \#0x1900}          & Wait for interrupt       \\
		\mnemonic{RETE}           & \texttt{COPI \#0x1A00}          & Return from exception    \\
		\mnemonic{RSET}           & \texttt{COPI \#0x1B00}          & System reset             \\
		\mnemonic{ETFR}           & \texttt{COPI \#0x1Cxy}          & Read exception state     \\
		\mnemonic{ETTR}           & \texttt{COPI \#0x1Dxy}          & Write exception state    \\
		\mnemonic{DMAR}           & \texttt{COPI \#0x28xx}          & DMA read to registers    \\
		\mnemonic{DMAW}           & \texttt{COPI \#0x29xx}          & DMA write from registers \\
		\mnemonic{DMAT}           & \texttt{COPI \#0x2Ann}          & DMA memory transfer      \\
		\mnemonic{MULU}           & \texttt{COPI \#0x3000}          & Unsigned multiply        \\
		\mnemonic{MULS}           & \texttt{COPI \#0x3100}          & Signed multiply          \\
		\mnemonic{DIVU}           & \texttt{COPI \#0x3200}          & Unsigned divide          \\
		\mnemonic{DIVS}           & \texttt{COPI \#0x3300}          & Signed divide            \\
		\bottomrule
	\end{tabular}
	\caption{Meta-Instructions}
	\label{tab:meta-instructions}
\end{table}

\section{Instruction Timing}

Ordinary processing-unit instructions use exactly 6 execution phases and complete in 6 clock cycles when no bus wait
states are required:

\begin{enumerate}
	\item Instruction Fetch (High Word)
	\item Instruction Fetch (Low Word)
	\item Decode and Register Fetch
	\item Execute and Address Calculation
	\item Memory Access
	\item Write Back
\end{enumerate}

Coprocessor calls (\mnemonic{COPI}, \mnemonic{COPR} and every meta-instruction that assembles to one) take a second
6-cycle slot for the coprocessor dispatch that begins at the next phase 0, for 12 cycles in total before the next
instruction is fetched.

This fixed sequence simplifies hardware design. Delayed bus acknowledgement may hold a phase and add wait cycles. DMA
meta-instructions dispatch through the normal 6-cycle \mnemonic{COPI} path and then add DMA command and transfer
cycles.

\section{Next Chapters}

The following chapters provide detailed documentation for each instruction group:

\begin{itemize}
	\item \textbf{Chapter~\ref{ch:alu-instructions}}: ALU instructions (arithmetic, logical, compare)
	\item \textbf{Chapter~\ref{ch:memory-instructions}}: Load and store operations
	\item \textbf{Chapter~\ref{ch:control-flow}}: Branches, jumps, and address operations
	\item \textbf{Chapter~\ref{ch:coprocessor}}: Coprocessor interface
	\item \textbf{Chapter~\ref{ch:meta-instructions}}: Meta-instruction cross-reference
\end{itemize}
